\documentclass[kvant.tex]{subfiles}

\begin{document}
\newline
5. а) Если \(p \neq 2\), то \(\displaystyle \alpha = \sum_{i=1}^{k} \left[ \frac{n}{p^i} \right]\), где k таково, что \(p^k\leq n < p^{k+1}\).\par
Если же p = 2, то \(\displaystyle \alpha = n + \sum_{i=1}^{k} \left[ \frac{n}{2^i} \right] \), где \(2^k \leq n < 2^{k+1}\).\par
б) \(\alpha = \left[ \frac{2n+1}{p} \right] - \left[ \frac{n}{p} \right] + \left[\frac{2n+1}{p^2} \right] + \left[\frac{n}{p^2} \right] + ... + \left[\frac{2n+1}{p^k} \right] - \left[\frac{n}{p^k} \right]\), где \(p^k \leq 2n+1 < p^{k+1}\).\par
{У к а з а н и е}. Воспользуйтесь тем, что \((2n+1)!! = \frac{(2n+1)!}{(2n)!!}\).\par
6) а) \(1\frac{1}{16}\); \(1\frac{3}{4}\); \(2\frac{7}{16}\); \(3\frac{1}{8}\); \(3\frac{3}{16}\);\par
б) \(-1 \leq x < 0\); \(\sqrt{8} \leq x < 3\); \(\sqrt{11} \leq x \leq \sqrt{14}\); \(4 \leq x < 17\).\par
в) \(\sqrt{8} \).\par
7. 30 \emph{км}.\par
8. [f(a)] + [f(a + 1)] + [f(a + 2)] + ... + [f(b)] + b - a + 1.\par
9. 137. У к а з а н и е. Координаты точек, лежащих в данном круге, удовлетворяют неравенству \(x^2+y^2 \geq 6,5^2\).
\end{document}
